\setcounter{section}{27}

  \zwischenueberschrift{Übungsaufgaben}

\inputaufgabe
{}
{

Es sei
\maabb {p} { E } { M
} {}
ein
\definitionsverweis {differenzierbares Vektorbündel}{}{}
über einer
\definitionsverweis {differenzierbaren Mannigfaltigkeit}{}{}
$M$, das mit einem
\definitionsverweis {linearen Zusammenhang}{}{}
$\nabla$ versehen sei. Es sei $L$ eine weitere differenzierbare Mannigfaltigkeit und
\maabb {\psi} { L } { M
} {}
eine
\definitionsverweis {stetig differenzierbare Abbildung}{}{}
mit dem
\definitionsverweis {zurückgezogenen Zusammenhang}{}{}
$\psi^*\nabla$ auf dem
\definitionsverweis {zurückgezogenen Vektorbündel}{}{}
$\psi^*E$ über $L$. Zeige, dass ein stetig differenzierbarer Schnitt
\maabb {s} { L } { E
} {}
genau dann
\definitionsverweis {horizontal}{}{}
bezüglich $\nabla$ ist, wenn der zugehörige Schnitt
\maabbdisp {s} { L } { \psi^*E
} {}
horizontal bezüglich  $\psi^* \nabla$ ist.

}
{} {}

\inputaufgabe
{}
{

Es sei
\maabb {p} {M \times \R^r } { M
} {}
das triviale
\definitionsverweis {Vektorbündel}{}{}
über einer
\definitionsverweis {differenzierbaren Mannigfaltigkeit}{}{}
$M$ versehen mit dem
\definitionsverweis {trivialen Zusammenhang}{}{}
$\nabla$. Es sei $L$ eine weitere differenzierbare Mannigfaltigkeit und
\maabb {\psi} { L } { M
} {}
eine
\definitionsverweis {stetig differenzierbare Abbildung}{}{.}
Zeige, dass der
\definitionsverweis {zurückgezogene Zusammenhang}{}{}
$\psi^*\nabla$
\zusatzklammer {auf $L \times \R^r$} {} {}
ebenfalls trivial ist.

}
{} {}

\inputaufgabe
{}
{

Es sei
\maabb {p} { E } { M
} {}
ein
\definitionsverweis {differenzierbares Vektorbündel}{}{}
über einer
\definitionsverweis {differenzierbaren Mannigfaltigkeit}{}{}
$M$, das mit einem
\definitionsverweis {linearen Zusammenhang}{}{}
$\nabla$ versehen sei. Es sei

\mavergleichskette
{\vergleichskette
{ I
}
{ \subseteq }{ \R
}
{  }{
}
{    }{
}
{   }{
}
}
{}{}{}
ein
\definitionsverweis {offenes Intervall}{}{}
und
\maabb {\gamma} { I } { M
} {}
eine
\definitionsverweis {stetig differenzierbare Kurve}{}{}
mit dem
\definitionsverweis {zurückgezogenen Zusammenhang}{}{}
$\gamma^*\nabla$ auf dem
\definitionsverweis {zurückgezogenen Vektorbündel}{}{}
$\gamma^*E$ über $I$. Bestimme die
\definitionsverweis {Christoffelsymbole}{}{}
von
\mathl{\gamma^* \nabla}{.}

}
{} {}

\inputaufgabegibtloesung
{}
{

Wir betrachten die trigonometrische Parametrisierung
\maabbeledisp {\psi} {\R } { S^1
} { t } { \left(  \cos  t   , \,   \sin  t                   \right)
} {}
des Einheitskreises. Es sei

\mavergleichskette
{\vergleichskette
{ c
}
{ \in }{ [0, 2 \pi[
}
{  }{
}
{    }{
}
{   }{
}
}
{}{}{}
fixiert. Auf dem trivialen Vektorbündel
\maabbdisp {} {S^1 \times \R^2} { S^1
} {}
sei ein
\definitionsverweis {linearer Zusammenhang}{}{}
$\nabla$ gegeben derart, dass der längs $\psi$
\definitionsverweis {zurückgezogene Zusammenhang}{}{}
$\psi^* \nabla$ durch die
\definitionsverweis {Christoffelsymbole}{}{}
\zusatzklammer {der Index für die einzige Ableitungsrichtung wird weggelassen} {} {}

\mavergleichskettedisp
{\vergleichskette
{ \begin{pmatrix} \tilde{\Gamma}_1^1 (t)   &  \tilde{\Gamma}_1^2 (t)   \\ \tilde{\Gamma}_2^1(t)  &  \tilde{\Gamma}_2^2(t)  \end{pmatrix}
}
{ =} { \begin{pmatrix}  0  &   - { \frac{  c  }{  2 \pi   } }   \\  { \frac{  c  }{  2 \pi  } }  &  0  \end{pmatrix}
}
{ } {
}
{ } {
}
{ } {
}
}
{}{}{}
gegeben ist. Wir betrachten zu einem Punkt

\mavergleichskette
{\vergleichskette
{ Q
}
{ = }{ \begin{pmatrix}  a  \\b   \end{pmatrix}
}
{ \in }{ \R^2
}
{    }{
}
{   }{
}
}
{}{}{}
die Abbildung
\maabbeledisp {\Psi_{Q}} {\R} { S^1 \times \R^2
} { t } { \left(  \cos  t   , \,   \sin  t  , \,   M(t)  \begin{pmatrix}  a  \\b   \end{pmatrix}                 \right)
} {}
mit

\mavergleichskettedisp
{\vergleichskette
{ M(t)
}
{ =} { \begin{pmatrix}  \cos  { \frac{  ct  }{  2 \pi   } }   &   -  \sin  { \frac{  ct  }{  2 \pi   } }    \\  \sin  { \frac{  ct  }{  2 \pi  } }   &  \cos  { \frac{  ct  }{  2 \pi   } }  \end{pmatrix}
}
{ } {
}
{ } {
}
{ } {
}
}
{}{}{.}
\aufzaehlungfuenf{Bestimme $\Psi_Q(0)$.
}{Bestimme $\Psi_Q(2 \pi)$.
}{Zeige, dass  $\Psi_Q$ eine
\definitionsverweis {horizontale Liftung}{}{}
längs $\psi$  ist.
}{Zeige, dass
\maabbeledisp {} {\R^2} { \R^2
} { Q } { \Psi_Q(2 \pi)
} {,}
eine
\definitionsverweis {lineare Isometrie}{}{}
ist
\zusatzklammer {der Basispunkt $(1,0)$ wird hier nicht aufgeführt} {} {.}
}{Zeige, dass
\maabbeledisp {} {\Z} { \operatorname{SL}_{  2  } \! \left(  \R  \right)
} { k } { \left(  Q \mapsto \Psi_Q(2k \pi )  \right)
} {,}
ein
\definitionsverweis {Gruppenhomomorphismus}{}{}
ist.
}

}
{} {}

\inputaufgabegibtloesung
{}
{

Bestimme für die Kurve
\maabbeledisp {\psi} {\R} { {\mathbb H}
} { t } { \left(  t   , \,  e^t                   \right)
} {,}
in die
\definitionsverweis {Halbebene von Poincaré}{}{}
${\mathbb H}$ die
\definitionsverweis {Christoffelsymbole}{}{}
für den
\definitionsverweis {zurückgezogenen Zusammenhang}{}{}
\zusatzklammer {zum
\definitionsverweis {Levi-Civita-Zusammenhang}{}{}
auf ${\mathbb H}$} {} {}
auf $\R$.

}
{} {}

Beachte in der folgenden Aufgabe, dass es nach
[[Topologische Räume/Relatives Produkt/Hintereinanderschaltung/Aufgabe|Aufgabe 11.23]]
egal ist, ob man den Rückzug des Vektorbündels nach $Y$ direkt oder über $Z$ bestimmt.

\inputaufgabe
{}
{

Es sei
\maabb {p} { E } { M
} {}
ein
\definitionsverweis {differenzierbares Vektorbündel}{}{}
über einer
\definitionsverweis {differenzierbaren Mannigfaltigkeit}{}{}
$M$, das mit einem
\definitionsverweis {linearen Zusammenhang}{}{}
$\nabla$ versehen sei. Es seien $Y,Z$ weitere differenzierbare Mannigfaltigkeiten mit
\definitionsverweis {stetig differenzierbare Abbildungen}{}{}

\mathdisp {Y \stackrel{\varphi}{\longrightarrow} Z \stackrel{\psi}{\longrightarrow}  M} { . }

Zeige, dass für die
\definitionsverweis {zurückgezogenen Zusammenhänge}{}{}
auf
\mathl{\left(  \psi \circ \varphi  \right)^* E}{} die Beziehung

\mavergleichskettedisp
{\vergleichskette
{ \left(  \psi \circ \varphi  \right)^*(\nabla)
}
{ =} { \varphi^*  \left(  \psi^*\nabla  \right)
}
{ } {
}
{ } {
}
{ } {
}
}
{}{}{}
gilt.

}
{} {}

\inputaufgabe
{}
{

Es sei $M$ eine
\definitionsverweis {riemannsche Mannigfaltigkeit}{}{}
versehen mit dem
\definitionsverweis {Levi-Civita-Zusammenhang}{}{}
und sei
\maabbdisp {\gamma} { I } { M
} {}
eine zweifach stetig
\definitionsverweis {differenzierbare Kurve}{}{.}
Zeige, dass $\gamma$ genau dann eine
\definitionsverweis {geodätische Kurve}{}{}
ist, wenn die Ableitungskurve
\maabb {\gamma'} { I } { TM
} {}
ein
\definitionsverweis {horizontaler Schnitt}{}{}
im
\definitionsverweis {Tangentialbündel}{}{}
$TM$ ist.

}
{} {}

\inputaufgabe
{}
{

Es sei

\mavergleichskette
{\vergleichskette
{ I
}
{ \subseteq }{ \R
}
{  }{
}
{    }{
}
{   }{
}
}
{}{}{}
ein
\definitionsverweis {offenes Intervall}{}{}
mit einer positiven stetig differenzierbaren Funktion
\maabbdisp {g} { I } { \R_+
} {,}
die wir als eine
\definitionsverweis {riemannsche Metrik}{}{}
auf $I$ auffassen. Charakterisiere die
\definitionsverweis {geodätischen Kurven}{}{}
\maabbdisp {\gamma} { J } { I
} {,}
wobei $J$ ein weiteres Intervall bezeichnet.

}
{} {}

\inputaufgabegibtloesung
{}
{

Bestimme für die Kurve
\maabbeledisp {\psi} {\R} { {\mathbb H}
} { t } { \left(  t   , \,  e^t                   \right)
} {,}
in die
\definitionsverweis {Halbebene von Poincaré}{}{}
${\mathbb H}$ die
\definitionsverweis {tangentiale Beschleunigung}{}{}
für jedes $t$.

}
{} {}

\inputaufgabe
{}
{

Man gebe ein Beispiel für eine
\definitionsverweis {zusammenhängende}{}{}
\definitionsverweis {riemannsche Mannigfaltigkeit}{}{}
$M$ und zwei Punkte

\mavergleichskette
{\vergleichskette
{ P,Q
}
{ \in }{ M
}
{  }{
}
{    }{
}
{   }{
}
}
{}{}{,}
die nicht durch eine
\definitionsverweis {geodätische Kurve}{}{}
verbunden werden können.

}
{} {}

\inputaufgabe
{}
{

Man gebe ein Beispiel für eine
\definitionsverweis {zusammenhängende}{}{}
\definitionsverweis {riemannsche Mannigfaltigkeit}{}{}
$M$ und zwei Punkte

\mavergleichskette
{\vergleichskette
{ P,Q
}
{ \in }{ M
}
{  }{
}
{    }{
}
{   }{
}
}
{}{}{,}
die durch mehrere
\definitionsverweis {geodätische Kurven}{}{}
miteinander verbunden werden können, die alle die gleiche Länge haben.

}
{} {}

\inputaufgabe
{}
{

Man gebe ein Beispiel für eine
\definitionsverweis {zusammenhängende}{}{}
\definitionsverweis {riemannsche Mannigfaltigkeit}{}{}
$M$ und zwei Punkte

\mavergleichskette
{\vergleichskette
{ P,Q
}
{ \in }{ M
}
{  }{
}
{    }{
}
{   }{
}
}
{}{}{,}
die durch mehrere
\definitionsverweis {geodätische Kurven}{}{}
miteinander verbunden werden können, die nicht die gleiche Länge haben.

}
{} {}

\inputaufgabegibtloesung
{}
{

Bestätige, dass das System von gewöhnlichen Differentialgleichungen zweiter Ordnung

\mavergleichskettedisp
{\vergleichskette
{ \gamma_1^{\prime \prime} (t)
}
{ =} { { \frac{   2  }{  \gamma_2(t) } } \gamma_1'(t) \gamma_2'(t)
}
{ } {
}
{ } {
}
{ } {
}
}
{}{}{}
und

\mavergleichskettedisp
{\vergleichskette
{ \gamma_2^{\prime \prime} (t)
}
{ =} { -  { \frac{   1  }{  \gamma_2(t) } }  \gamma_1'(t) \gamma_1'(t) + { \frac{   1  }{  \gamma_2(t) } }    \gamma_2'(t) \gamma_2'(t)
}
{ } {
}
{ } {
}
{ } {
}
}
{}{}{}
durch die Kurven

\mavergleichskettedisp
{\vergleichskette
{ \gamma(t)
}
{ =} { \begin{pmatrix}  s+r \tanh  t  \\ { \frac{   r  }{  \cosh  t  } }    \end{pmatrix}
}
{ } {
}
{ } {
}
{ } {
}
}
{}{}{}
zu

\mavergleichskette
{\vergleichskette
{ s
}
{ \in }{ \R
}
{  }{
}
{    }{
}
{   }{
}
}
{}{}{,}

\mavergleichskette
{\vergleichskette
{ r
}
{ \in }{ \R_+
}
{  }{
}
{    }{
}
{   }{
}
}
{}{}{}
gelöst wird.

}
{} {}

\inputaufgabe
{}
{

Wir betrachten Kreise mit Mittelpunkt
\mathl{(s,0)}{} und Radius

\mavergleichskette
{\vergleichskette
{ r
}
{ > }{ 0
}
{  }{
}
{    }{
}
{   }{
}
}
{}{}{.}
Charakterisiere, wann sich zwei solche Kreise schneiden.

}
{} {}

\inputaufgabe
{}
{

Es sei

\mavergleichskette
{\vergleichskette
{ P
}
{ = }{(a,b)
}
{ \in }{ {\mathbb H}
}
{    }{
}
{   }{
}
}
{}{}{}
ein Punkt in der
\definitionsverweis {Halbebene von Poincaré}{}{}
und sei

\mavergleichskette
{\vergleichskette
{ v
}
{ = }{ (v_1,v_2)
}
{ \in }{ \R^2
}
{    }{
}
{   }{
}
}
{}{}{}
ein
\definitionsverweis {Tangentenvektor}{}{}
an $P$.
\aufzaehlungzwei {Bestimme einen Halbkreis mit Mittelpunkt auf der $x$-Achse
\zusatzklammer {mit vertikalen Geraden als Extremfall} {} {,}
der durch $P$ verläuft und tangential zu $\R v$ ist.
} {Bestimme eine
\definitionsverweis {geodätische Kurve}{}{,}
die diesen Tangentenvektor realisiert.
}

}
{} {}

\inputaufgabe
{}
{

Es sei
\maabb {\psi} { L } { M
} {}
eine
\definitionsverweis {Isometrie}{}{}
zwischen
\definitionsverweis {riemannschen Mannigfaltigkeiten}{}{}
\mathkor {} {L} {und} {M} {.}
\aufzaehlungdrei{Zeige, dass der
\definitionsverweis {zurückgezogene Zusammenhang}{}{}
zum
\definitionsverweis {Levi-Civita-Zusammenhang}{}{}
auf $M$ gleich dem Levi-Civita-Zusammenhang auf $L$ ist.
}{Zeige, dass sich die
\definitionsverweis {tangentialen Beschleunigungen}{}{}
zu einer zweifach differenzierbaren Kurve
\maabb {\gamma} { I } { L
} {}
bzw.
\maabb {\psi \circ \gamma} { I } { M
} {}
entsprechen.
}{Zeige, dass sich die
\definitionsverweis {geodätischen Kurven}{}{}
in $L$ und in $M$ entsprechen.
}

}
{} {}

Es sei $M$ eine
\definitionsverweis {riemannsche Mannigfaltigkeit}{}{,}
wobei das
\definitionsverweis {Tangentialbündel}{}{}
$TM$ mit dem
\definitionsverweis {Levi-Civita-Zusammenhang}{}{}
versehen sei. Es sei
\maabb {\gamma} { I } { M
} {}
eine
\definitionsverweis {differenzierbare Kurve}{}{.}
Man sagt, dass ein längs $\gamma$ definiertes differenzierbares
\definitionsverweis {Vektorfeld}{}{}
\maabb {F} { I } { TM
} {}
\definitionswort {parallel längs}{}
$\gamma$ ist, wenn

\mavergleichskettedisp
{\vergleichskette
{ (\nabla_{\partial} F)(t)
}
{ =} { 0
}
{ } {
}
{ } {
}
{ } {
}
}
{}{}{}
für alle

\mavergleichskette
{\vergleichskette
{ t
}
{ \in }{ I
}
{  }{
}
{    }{
}
{   }{
}
}
{}{}{}
gilt.

Für die folgende Aufgabe vergleiche
[[Differenzierbare Hyperfläche/Zusammenhang/Kurve/Parallel und horizontal/Aufgabe|Aufgabe 25.13]].

\inputaufgabe
{}
{

Es sei

\mavergleichskette
{\vergleichskette
{ Y
}
{ \subseteq }{ W
}
{ \subseteq }{ \R^n
}
{    }{
}
{   }{
}
}
{}{}{,}
$W$
\definitionsverweis {offen}{}{,}
eine
\definitionsverweis {differenzierbare Hyperfläche}{}{,}
die wir auch also eine
\definitionsverweis {riemannsche Mannigfaltigkeit}{}{}
über den umgebenden Raum $\R^n$ mit dem
\definitionsverweis {Standardskalarprodukt}{}{}
auffassen.
Es sei
\maabb {\gamma} { I } { Y
} {}
eine
\definitionsverweis {differenzierbare Kurve}{}{}
und es sei
\maabb {F} { I } { TY \subseteq Y \times \R^n
} {}
ein differenzierbares
\definitionsverweis {Vektorfeld}{}{}
längs $\gamma$.  Zeige, dass $F$ genau dann
\definitionswort {parallel längs}{}
$\gamma$ im Sinne von
[[Hyperfläche/Kurve/Paralleles Vektorfeld/Definition|Definition 6.6]]
ist, wenn $F$
\definitionsverweis {parallel}{}{}
längs $\gamma$ bezüglich des
\definitionsverweis {Levi-Civita-Zusammenhangs}{}{}
auf $TY$ ist.

}
{} {}

\inputaufgabe
{}
{

Es sei $M$ eine
\definitionsverweis {riemannsche Mannigfaltigkeit}{}{,}
wobei das
\definitionsverweis {Tangentialbündel}{}{}
$TM$ mit dem
\definitionsverweis {Levi-Civita-Zusammenhang}{}{}
versehen sei. Es sei
\maabb {\gamma} { I } { M
} {}
eine zweifach
\definitionsverweis {differenzierbare Kurve}{}{.}
Zeige, dass $\gamma$ genau dann eine
\definitionsverweis {geodätische Kurve}{}{}
ist, wenn die Ableitung
\maabbdisp {\gamma'} { I } { TM
} {}
ein
\definitionsverweis {paralleles Vektorfeld}{}{}
ist.

}
{} {}

\inputaufgabe
{}
{

Zeige, dass auf dem
\definitionsverweis {euklidischen Raum}{}{}
$\R^n$ die
\definitionsverweis {riemannsche Metrik}{}{}
gleich dem
\definitionsverweis {euklidischen Abstand}{}{}
ist.

}
{} {}

  \zwischenueberschrift{Aufgaben zum Abgeben}

\inputaufgabe
{4}
{

Wir betrachten $\R_+$ mit der durch

\mavergleichskette
{\vergleichskette
{ g(t)
}
{ = }{ t^2
}
{  }{
}
{    }{
}
{   }{
}
}
{}{}{}
gegebenen
\definitionsverweis {riemannschen Metrik}{}{.}
Bestimme die
\definitionsverweis {geodätischen Kurven}{}{}
in $\R_+$.

}
{} {}

\inputaufgabe
{6 (3+3)}
{

Wir betrachten die Kurve
\maabbeledisp {\psi} {\R} { {\mathbb H}
} { t } { \left(  t   , \,   1+t^2                   \right)
} {,}
in die
\definitionsverweis {Halbebene von Poincaré}{}{}
${\mathbb H}$.
\aufzaehlungzwei {Bestimme die
\definitionsverweis {Christoffelsymbole}{}{}
für den
\definitionsverweis {zurückgezogenen Zusammenhang}{}{}
\zusatzklammer {zum
\definitionsverweis {Levi-Civita-Zusammenhang}{}{}
auf ${\mathbb H}$} {} {}
auf $\R$.
} {Bestimme die
\definitionsverweis {tangentiale Beschleunigung}{}{}
von $\psi$ für jedes $t$.
}

}
{} {}

\inputaufgabe
{4}
{

Bestimme die
\definitionsverweis {geodätischen Kurven}{}{}
auf der
\definitionsverweis {hyperbolischen Kreisscheibe}{}{.}

}
{} {}

\inputaufgabe
{4}
{

Es sei

\mavergleichskette
{\vergleichskette
{ n
}
{ \geq }{ 2
}
{  }{
}
{    }{
}
{   }{
}
}
{}{}{}
und sei

\mavergleichskettedisp
{\vergleichskette
{ M
}
{ =} { \R^n \setminus \{P_1  , \ldots , P_m \}
}
{ \subseteq} { \R^n
}
{ } {
}
{ } {
}
}
{}{}{,}
aufgefasst als
\definitionsverweis {riemannsche Mannigfaltigkeit}{}{}
mit dem induzierten
\definitionsverweis {Standardskalarprodukt}{}{.}
Zeige, dass auf $M$ die
\definitionsverweis {riemannsche Metrik}{}{}
gleich dem
\definitionsverweis {euklidischen Abstand}{}{}
ist.

}
{} {}

 [[Kategorie:Latexseite]]
